% Timing diagram of the CPU's interrupt request interface: one hardware
% interrupt taken at the boundary it arrives at, followed by a second request
% that arrives while the service routine is still running and is therefore made
% to wait until the boundary after the RTI.
%
% The interface is AXI-stream shaped -- irq_valid_i / irq_ready_o / irq_addr_i --
% and deliberately knows nothing about daisy chaining, which is a matter for the
% adaptation layer around the CPU.  See src/interrupt/README.md.
%
% All values below are taken from a GHDL simulation of
% test/prog_int_waveform.asm at the default (zero) memory latencies: t=0 is the
% clock cycle from 270 ns to 280 ns.  Changing that program, or the number of
% cycles any instruction in it takes, invalidates them.  The addresses are:
%
%   000C-000E  padding "MOVE R2, R2" in the main program
%   0026       ISR1, which requests the second interrupt
%   002A-002D  the last three instructions of ISR1, ending in RTI at 002D
%   002E       ISR2

\documentclass{standalone}
\usepackage{timing}

% Four-digit addresses do not fit in a default-width box.
\renewcommand{\wavescale}{1.15}

\begin{document}
% Drawing convention, the same one src/cpu_main/timing.tex uses: handshake bits
% are always drawn with their true value; a bus is greyed out ("x") in every
% cycle in which it carries nothing this diagram is about -- irq_addr_i while
% irq_valid_i is low (test/interrupt.vhd drives all-ones there on purpose),
% prep_stage_i.addr and resume_pc except in a cycle an instruction retires in or
% a request is taken in, and fetch_addr_o except when fetch_valid_o is high.
%
% There is no interrupt module between the pins and WRITE, and nothing in the
% handshake is registered: WRITE reads irq_valid_i and irq_addr_i in the cycle an
% instruction retires, and irq_ready_o, fetch_valid_o, and fetch_addr_o are all
% combinational in them.  The only registers are irq_active and irq_r15 (and
% irq_r14, not drawn), which change on the edge that ends a cycle with a
% transfer or an RTI.
%
% t=1   Request 1 arrives in the cycle 000D retires, and is taken in that same
%       cycle: accepted, redirected to 0026, and resume_pc (000E) saved.
% t=11  Request 2 arrives inside ISR1.  It is refused at 002A, 002C, and at the
%       RTI at 002D (t=13), because irq_active is still high in that cycle.
% t=18  The first instruction back at the return address, 000E, retires, and
%       request 2 is taken, saving 000F.
\begin{wave}{10}{20}
 \nextwave{irq\_valid\_i}                 \bit{0}{1} \bit{1}{1} \bit{0}{9} \bit{1}{8} \bit{0}{1}
 \nextwave{irq\_addr\_i}                  \unknown[x]{1} \known{0026}{1} \unknown[x]{9} \known{002E}{8} \unknown[x]{1}
 \nextwave{irq\_ready\_o}                 \bit{0}{1} \bit{1}{1} \bit{0}{16} \bit{1}{1} \bit{0}{1}
 \nextwave{WRITE/inst\_done\_o}           \bit{1}{2} \bit{0}{5} \bit{1}{1} \bit{0}{1} \bit{1}{1} \bit{0}{1} \bit{1}{3} \bit{0}{4} \bit{1}{1} \bit{0}{1}
 \nextwave{WRITE/prep\_stage\_i.addr}     \known{000C}{1} \known{000D}{1} \unknown[x]{5} \known{0026}{1} \unknown[x]{1} \known{0028}{1} \unknown[x]{1} \known{002A}{1} \known{002C}{1} \known{002D}{1} \unknown[x]{4} \known{000E}{1} \unknown[x]{1}
 \nextwave{WRITE/resume\_pc}              \unknown[x]{1} \known{000E}{1} \unknown[x]{16} \known{000F}{1} \unknown[x]{1}
 \nextwave{WRITE/irq\_active}             \bit{0}{2} \bit{1}{12} \bit{0}{5} \bit{1}{1}
 \nextwave{WRITE/irq\_r15}                \unknown[x]{2} \known{000E}{17} \known{000F}{1}
 \nextwave{WRITE/fetch\_valid\_o}         \bit{0}{1} \bit{1}{1} \bit{0}{11} \bit{1}{1} \bit{0}{4} \bit{1}{1} \bit{0}{1}
 \nextwave{WRITE/fetch\_addr\_o}          \unknown[x]{1} \known{0026}{1} \unknown[x]{11} \known{000E}{1} \unknown[x]{4} \known{002E}{1} \unknown[x]{1}
\end{wave}

\end{document}
